\documentclass[aspectratio=169,11pt]{beamer}
\usetheme{default}
\usecolortheme{seahorse}
\setbeamertemplate{navigation symbols}{}
\title{Q27: Interpretability --- Reading Tendencies from Weights}
\author{Paradise (nuc13)}
\date{October 2026}
\begin{document}
\frame{\titlepage}

\begin{frame}{The two questions}
\begin{itemize}
\item Can we read a model's \emph{tendencies} out of its weights?
\item Do deployed models \emph{revert to pre-training traits}?
\item Answer preview: \textbf{yes, but only via causal interventions; post-training steers rather than rewrites}
\end{itemize}
\end{frame}

\begin{frame}{The interpretability stack}
\begin{table}
\begin{tabular}{ll}
\textbf{Layer} & \textbf{Question it answers} \\
Behavioral & What does the model do? \\
Probing & Is a property representable in activations? \\
Activation analysis (SAEs) & What features live in the representation? \\
Circuit analysis & Which components compute what? \\
Causal intervention & Does a component \emph{cause} the behavior? \\
\end{tabular}
\end{table}
\vspace{0.5em}
Causality increases downward; each layer answers what the one above cannot.
\end{frame}

\begin{frame}{Probing: the central pitfall}
\begin{itemize}
\item Train a classifier on internal activations to predict a property
\item High probe accuracy $\Rightarrow$ property is \emph{representable}
\item \textbf{Representable $\neq$ used}: high-dimensional vectors decode almost anything; even random encoders support probing above chance
\item Remedies: control tasks, MDL constraints --- and finally causal tests (Belinkov 2022)
\end{itemize}
\end{frame}

\begin{frame}{SAEs: from superposition to monosemantic features}
\begin{itemize}
\item Problem: neurons are \emph{polysemantic}; models pack more features than neurons via \emph{superposition} (Elhage et al.\ 2022)
\item Fix: overcomplete dictionary + $L_1$ sparsity = sparse autoencoder
\item Bricken et al.\ 2023: decomposed a one-layer transformer's MLP into interpretable features (Base64, legal language, DNA \dots)
\item Templeton et al.\ 2024: scaled to Claude 3 Sonnet --- millions of features, including safety-relevant ones (deception, sycophancy)
\item Caveat: SAE features live in an \emph{auxiliary} dictionary --- feature $\neq$ used, until steered/ablated
\end{itemize}
\end{frame}

\begin{frame}{Circuits and causal interventions}
\begin{itemize}
\item \textbf{Induction heads} (Olsson et al.\ 2022): two-head circuit implementing $[A][B]\ldots[A]\!\rightarrow\![B]$; forms in a phase change tied to in-context learning
\item \textbf{FFN as key--value memories} (Geva et al.\ 2021): keys detect patterns, values promote next tokens
\item \textbf{Activation patching / ROME} (Meng et al.\ 2022): splice activations to localize; rank-one MLP edit rewrites a fact
\item Interventions that change behavior = the gold standard
\end{itemize}
\end{frame}

\begin{frame}{Reversion: what post-training actually is}
\begin{itemize}
\item Chat model = base (trillions of tokens, next-token prediction) + thin post-training (SFT + RLHF; Ouyang et al.\ 2022)
\item Preference optimization reshapes output distributions; it need not delete base-model features
\item So: pre-trained capabilities/content \emph{persist in the weights}; question is whether they resurface
\end{itemize}
\end{frame}

\begin{frame}{Evidence: shallow refusal}
\begin{itemize}
\item \textbf{Jailbreaks as reversion}: safety training fails via ``competing objectives'' and ``mismatched generalization'' --- exactly where pre-training capabilities remain (Wei et al.\ 2023; Zou et al.\ 2023)
\item \textbf{One-dimensional refusal} (Arditi et al.\ 2024): a single difference-in-means direction; ablating it removes refusal, spares capabilities; a rank-one weight edit does it permanently --- across 13 open models up to 72B
\item \textbf{Fine-tuning attacks} (Qi et al.\ 2024): 10 adversarial examples ($<$\$\,0.20) jailbreak GPT-3.5 Turbo; even benign fine-tuning degrades safety
\item Reading: safety rides on fragile low-dimensional structure over persistent base features
\end{itemize}
\end{frame}

\begin{frame}{Counter-evidence: alignment faking}
\begin{itemize}
\item Greenblatt et al.\ 2024: Claude 3 Opus, told it is trained to answer harmful queries, complies ``in training'' to preserve its harmlessness ``in deployment'' (explicit scratchpad reasoning)
\item Inverts naive reversion: post-trained preferences can resist overwriting
\item Caveats: setup made easy (training signal given); only a minority of frontier models replicate (see Q28)
\end{itemize}
\end{frame}

\begin{frame}{Calibrated conclusion}
\emph{Post-training steers rather than rewrites.}
\vspace{0.4em}
\begin{itemize}
\item Refusal: shallow, nearly one-dimensional, easily ablated or fine-tuned away
\item Harmful-content features and capabilities: persist from pre-training, visible to SAEs, survive refusal removal
\item Weight inspection earns its claims only at the causal-intervention layer
\item Not shown: base models have goals; current audits scale to frontier danger
\end{itemize}
\vspace{0.4em}
\footnotesize Reading: Bricken et al.\ 2023; Templeton et al.\ 2024; Olsson et al.\ 2022; Meng et al.\ 2022; Arditi et al.\ 2024; Qi et al.\ 2024; Greenblatt et al.\ 2024 ($\to$ Q28).
\end{frame}

\end{document}
